\section{Milestones and Timeline}
\label{sec:timeline}

Although the focus of each thrust is sequenced, we plan to operate around a
working end-to-end prototype from the first quarter, so that every component
is exercised in context throughout the project. The targets below refer to
the metrics of Section~\ref{sec:evaluation}, and each year ends with a
go/no-go review.

{\footnotesize\setlength{\tabcolsep}{4pt}\renewcommand{\arraystretch}{0.93}
\begin{longtable}{@{}p{0.06\linewidth}p{0.48\linewidth}p{0.385\linewidth}@{}}
\toprule
\textbf{Qtr} & \textbf{Deliverable} & \textbf{Success criterion}\\
\midrule
Q1--Q2 & Knowledge-base schema with provenance and confidence;
     documentation-ingestion agents; two reference networks (10--20 hosts) as
     infrastructure-as-code; Observer Agents and passive observation;
     service-dependency capture; first ZMap safety extensions
     (impact-budgeted scheduling, interlocks).
   & Schema expresses all reference-network ground-truth entities; $\geq 85\%$
     of hosts and services recovered; no measurable production impact under
     interlocks.\\
\addlinespace[2pt]
Q3--Q4 & Observation planning as constrained optimization; model coverage
     metric; knowledge-graph-native ZMap output and safe fingerprinting;
     first-generation Replication Agents (equivalence-class abstraction,
     interaction patterns, synthetic data with authorization-structure
     preservation). \textbf{First upstream ZMap contribution; end-to-end
     prototype.}
   & $\geq 95\%$ hosts/services and $\geq 90\%$ dependencies; coverage estimate
     within 10 points; 20-host twin built unattended at $\geq 5\times$ resource
     reduction. \textbf{Go/no-go:} end-to-end run with no manual intervention
     between stages.\\
\addlinespace[2pt]
Q5--Q6 & Logic-based invariants and attack-graph preservation; quantitative
     fidelity measure; mock synthesis with declared unsoundness regions;
     component-level Red Team Agents; \textbf{nginx hardening campaign}
     (memory-safety analysis, configuration-defect checker, differential
     proxy-upstream smuggling tests).
   & $\geq 90\%$ attack-path recall and $\geq 75\%$ precision; mock-backed surface
     correctly identified; $\geq 85\%$ of seeded single-component defects;
     \textbf{first confirmed defect reported upstream}.\\
\addlinespace[2pt]
Q7--Q8 & Multi-step planning (ChainReactor and UTE extended to networks;
     equivalence-class lifting); analyst verification workflow;
     continuous-operation cadences and cost reduction; evaluation on partner
     and campus networks; public release of framework, benchmark, and twinCTF
     challenges.
   & $\geq 70\%$ of seeded multi-step chains; planner scales to 100 hosts;
     $\leq 20\%$ of findings fail to reproduce; $\geq 10$ upstream reports
     total; artifacts reproducible from published commits.\\
\bottomrule
\end{longtable}}

The ZMap extensions gate the quality of the capture that the
first-generation Replication Agents consume, so they are scheduled across
Q1--Q4, with acceptance criteria stated in terms of downstream usability
instead of feature completeness. The fidelity measure gates the interpretation
of every result of Thrust~3, and it therefore precedes the analysis milestones
in Q5--Q6. Should that milestone miss its recall target, the remaining
quarters will proceed on the reference networks, where ground truth
substitutes for the measure, while the work on fidelity continues in parallel.
